% OWNER: sections/s2_5_orders.tex. Edit only this file.
\seclabel
\claim{\secnum{2.5}{sec:orders}\quad Orders are placed by REST and confirmed in the feed; \code{wei} sets their size}

\lead{The server accepts an order placed before the open.} An \code{AddOrder} request sent at 08:30:05 with a body taken from the web client's code named the date (\code{d}), the portfolio (\code{pfId}), the ticker (\code{prodId} = \code{A079550}, LIG Nex1), the side (\code{ls} = L), the entry code (\code{ex} = E) and the order's size field (\code{wei} = 0.3). Two more groups of fields set how the order trades. The price-type fields choose a limit at a fixed price (\code{limitPrc}), a limit $n$ ticks from the market (\code{limitIdx}) or a stop price (\code{stopPrc}); the timing fields \code{hm0} and \code{hm1} either send the order at once or spread its fills over a window, here 09:00 to 09:10. The server answered HTTP 200 with \texttt{\{"data":"","warnings":[]\}}.

\lead{The reply carries no order ID, so the wrapper finds its own order in the live feed.} The feed reply right after the order had \code{null} arguments (Section~\ref{sec:feed}); by 08:33 the feed's order list (\code{ords}) held a new record, ID 1649882. The wrapper matches a new record by ticker, side and size among the IDs that were absent before the write, and it holds one lock across the write and the search, so 2 identical orders cannot bind the same ID.

\lead{A lost reply leaves the wrapper unsure whether the order exists, so the wrapper never resends an order automatically (\figref{orderstates}).} After a timeout, a 5xx or a connection lost after the request left, the wrapper looks for the order in the feed for 15\,s, then in the REST \code{Portfolio/Orders} records created after the write; if neither holds it, the wrapper raises \code{AmbiguousState} and leaves the check to the user. A resend could place the order twice, and each extra order uses one of the 80 a day the wrapper allows. The one exception is a connection failure that provably never reached the server (\code{ConnectError}, \code{ConnectTimeout} or \code{PoolTimeout}): the wrapper sends that request once more. A 401 is renewed and sent once more, because ASP.NET refuses it before the order handler runs (Section~\ref{sec:auth}); a 400 is never retried.

% OWNER: figs/orderstates.tex. Edit only this file.
\fig[wide]{orderstates}{An order the server may have received is never resent; the server's order ends in \code{Done} whether filled or stopped}[Top: what one \code{AddOrder} call does after each kind of reply. Bottom: the order's \code{state} on the server (names from the web client's code)]{%
\begin{tikzpicture}[x=1mm,y=1mm,
  f/.style={font=\fontsize{7}{8.4}\selectfont},
  st/.style={rectangle,rounded corners=1.6mm,draw=navy,line width=0.6pt,f,text=navy,align=center,minimum width=24mm,minimum height=8.5mm,inner sep=1pt,fill=white},
  sv/.style={st,rounded corners=4.2mm,minimum width=22mm},
  un/.style={sv,draw=greyD,dashed,text=greyD},
  term/.style={st,fill=navy,text=white},
  tr/.style={-{Stealth[length=1.6mm]},line width=0.6pt,navy},
  trU/.style={-{Stealth[length=1.6mm]},line width=0.5pt,greyD,dashed},
  lb/.style={f,text=navy,inner sep=1pt,align=center},
  lbU/.style={f,text=greyD,inner sep=1pt,align=center},
  band/.style={font=\fontsize{7.5}{9}\selectfont\bfseries,text=navy,anchor=west}]
% ------------------------------------------------ wrapper band
\node[band] at (-22,33) {Wrapper: one \code{AddOrder} call};
\node[st] (S) at (32,0) {Sending\\\code{AddOrder}};
\node[st] (F) at (84,0) {Looking for the\\order in the feed};
\node[st] (Q) at (134,0) {Asking REST\\\code{Orders}};
\node[term] (T4) at (32,25) {\code{RequestError}\\raised};
\node[term] (TA) at (134,25) {\code{AmbiguousState}:\\order not resent};
\node[term] (TI) at (116,-17) {Order ID returned};
% self loops on S
\draw[tr] ([yshift=2mm]S.west) to[out=160,in=110,looseness=6] ([xshift=-7mm]S.north);
\node[lb,anchor=east,align=right] at (13,9) {401: renew the token,\\send once more};
\draw[tr] ([yshift=-2mm]S.west) to[out=-160,in=-110,looseness=6] ([xshift=-7mm]S.south);
\node[lb,anchor=east,align=right] at (13,-9) {connection never\\established: send once\\more; if that fails too,\\look for the order};
% S -> T4, S -> F
\draw[tr] (S.north) -- node[lb,right]{400: never retried} (T4.south);
\draw[tr] (S.east) -- node[lb,above]{200 with no\\order ID} node[lb,below,align=center]{or reply lost:\\timeout, 5xx or\\connection lost\\after sending} (F.west);
% F -> Q, F -> TI, Q -> TA, Q -> TI
\draw[tr] (F.east) -- node[lb,above]{not found} node[lb,below]{in 15\,s} (Q.west);
\draw[tr] (Q.north) -- node[lb,right,align=left]{not among records\\newer than the write} (TA.south);
\draw[tr] (F.south) |- (TI.west);
\node[lb,anchor=north] at (92,-18.2) {found: same ticker,\\\code{ex} and size};
\draw[tr] (Q.south) |- node[lb,pos=0.25,right,xshift=0.8mm,align=left]{found} (TI.east);
% ------------------------------------------------ separator
\draw[greyM,line width=0.4pt] (-22,-27) -- (158,-27);
\node[band] at (-22,-31) {Server: the order's \code{state}};
% ------------------------------------------------ server band
\node[un] (G) at (-10,-44) {\code{Generated}};
\node[sv] (A) at (28,-44) {\code{Accepted}};
\node[sv] (W) at (70,-44) {\code{Working}};
\node[sv] (D) at (112,-44) {\code{Done}};
\node[un] (X) at (9,-62) {\code{Deleted}};
\draw[trU] (G) -- (A);
\draw[trU] (A) -- node[lbU,above,yshift=0.8mm]{not observed} (W);
\draw[trU] (G) -- (X);
\draw[trU] (A) -- (X);
\node[lbU,anchor=east,align=right] at (-3,-58) {\code{DeleteOrder}};
\draw[tr] (W.north east) to[out=25,in=155] node[lb,above]{last share filled} (D.north west);
\draw[tr] (W.east) -- node[lb,below,align=center]{stopped} (D.west);
\draw[tr] (A.south) to[out=-35,in=-145] node[lb,below,align=center]{stopped at once (order 1650719)} (D.south);
\node[lb,anchor=west,align=left] at (125,-44) {filled if \code{cnclT} null;\\stopped if \code{cnclT} set;\\partly filled, then\\stopped if also \code{cumQty} $>$ 0};
% ------------------------------------------------ legend
\begin{scope}[shift={(-22,-78)}]
\node[st,minimum width=8mm,minimum height=4.5mm] at (4,0) {}; \node[f,anchor=west] at (9,0) {wrapper step};
\node[term,minimum width=8mm,minimum height=4.5mm] at (32,0) {}; \node[f,anchor=west] at (37,0) {end of the call};
\node[sv,minimum width=8mm,minimum height=4.5mm,rounded corners=2.2mm] at (63,0) {}; \node[f,anchor=west] at (68,0) {server state, observed};
\node[un,minimum width=8mm,minimum height=4.5mm,rounded corners=2.2mm] at (101,0) {}; \node[f,anchor=west] at (106,0) {in client code, not observed};
\draw[tr] (145,1.6) -- (152,1.6); \node[f,anchor=west] at (153,1.6) {rule or observed};
\draw[trU] (145,-1.8) -- (152,-1.8); \node[f,anchor=west] at (153,-1.8) {not observed};
\end{scope}
\end{tikzpicture}}{The wrapper never sends an order a second time when it cannot tell whether the first reached the server. It sends once more only in 2 cases that cannot create a duplicate: a connection that was never established (\code{ConnectError}, \code{ConnectTimeout}, \code{PoolTimeout}), and a 401, which ASP.NET returns before the request reaches its handler (D36). If the resend after a connection failure also fails, the order goes to the feed and REST search; there is no third send. A 403 or 429 on any request raises \code{BlockedError} and halts all writes (not drawn). One lock covers sending and finding, and IDs already returned are skipped, so 2 identical orders cannot be matched to the same order ID. Server states: \code{Accepted} was seen right after placing order 1650719, \code{Working} on the 08:30 order; \code{Generated} and \code{Deleted} appear only in the client code. Every finished order shows \code{Done}, so \code{cnclT} and \code{cumQty} tell a fill from a stop.}{\code{pytimefolio} \code{orders} and \code{transport} modules; live feed and \code{Orders} replies, 24 Nov 2026; DART Quant Team}


\lead{\code{wei} on an order is the order's size in percent of portfolio value; a target weight would have given a different share count.} Before the order the portfolio held 3 shares of LIG Nex1, 0.13\% of its value, and one share weighed about 0.043\%. Read as a target, \code{wei} = 0.3 would ask for about 4 more shares; read as a size, it is worth 7.1 shares. The order reported \code{qty2ord} = 7, so the wrapper converts a target into an order as target minus current weight. It rounds \code{wei} once to 0.01\% when it builds the order and sends and matches that same value; the server stored order sizes exactly as sent (0.04, 0.1 and 0.3 were observed). The server rounds \code{qty2ord} down to whole shares: a test that separates floor from round-to-nearest is in Section~\ref{sec:oneshare}.

\lead{The order was \code{Working} before the open and filled in batches of 1 to 2 shares about every 2 minutes, 7 shares over 7\,min 37\,s (\figref{fills}).} The feed, polled every 3\,s, showed 1 share filled at 09:00:30, 3 at 09:02:09, 4 at 09:04:10, 6 at 09:06:10 and 7 at 09:08:07, when the state changed to \code{Done}, 1\,min 53\,s before the window closed. The position ended at 10 shares, 0.43\% of the portfolio.

% OWNER: figs/fills.tex. Edit only this file.
\fig{fills}{Seven shares filled in batches of 1 or 2 about every 2\,min, ending at 09:08:07}[Shares filled (\code{cumQty}) of order 1649882 (LIG Nex1, \code{wei} 0.3), read from the live feed on 24 Nov 2026]{%
\begin{tikzpicture}[x=0.15mm,y=4.6mm,font=\fontsize{7}{8.4}\selectfont]
% x: seconds after 09:00 (0.15 mm per s); y: shares filled
\fill[navyXL] (0,0) rectangle (600,7.6);
% grid and axes
\foreach \y in {1,...,7} \draw[greyM,line width=0.2pt] (-60,\y) -- (660,\y);
\draw[greyD,line width=0.4pt] (-60,0) -- (660,0);
\draw[greyD,line width=0.4pt] (-60,0) -- (-60,7.6);
\foreach \y in {0,...,7} \node[anchor=east,text=greyD] at (-60,\y) {\y};
\foreach \m/\l in {0/09:00,120/09:02,240/09:04,360/09:06,480/09:08,600/09:10}
  {\draw[greyD] (\m,0) -- (\m,-0.18); \node[anchor=north,text=greyD] at (\m,-0.2) {\l};}
\node[anchor=north,text=greyD] at (300,-1.05) {Time (KST)};
\node[rotate=90,anchor=south,text=greyD] at (-128,3.5) {Shares filled of 7};
\node[text=navy,align=center,fill=navyXL,inner sep=1pt] at (470,1.5) {order window\\09:00 to 09:10};
% step line
\draw[navy,line width=1pt] (-60,0) -- (30,0) -- (30,1) -- (129,1) -- (129,3) -- (250,3) -- (250,4) -- (370,4) -- (370,6) -- (487,6) -- (487,7) -- (660,7);
% fill times
\foreach \t/\n/\l in {30/1/09:00:30,129/3/09:02:09,250/4/09:04:10,370/6/09:06:10}
  {\fill[navy] (\t,\n) circle (0.75mm); \node[anchor=south east,text=navy,inner sep=1.5pt] at (\t,\n) {\l};}
\fill[navy] (487,7) circle (0.75mm);
\node[anchor=north west,text=navy,inner sep=2pt,font=\fontsize{7}{8.4}\selectfont\bfseries] at (487,7) {09:08:07 \code{Done}};
% position weight after each fill
\node[anchor=east,text=greyD] at (-60,8.35) {Position weight (\%)};
\foreach \t/\w in {30/0.17,129/0.26,250/0.30,370/0.38,487/0.43}
  {\node[text=greyD] at (\t,8.35) {\w};}
\end{tikzpicture}}{Each dot is a fill time read from the live feed, which was polled every 3\,s, so each time is accurate to $\pm$3\,s. Shading marks the order's window (\code{hm0} 09:00 to \code{hm1} 09:10); the order was \code{Working} from 08:33 and changed to \code{Done} with the last share, 7 min 37\,s after the first. The top row gives the position's weight in \% of net assets after each fill, counting the 3 shares held before the order. Average fill price KRW~352,714.29.}{Monitor log of the live feed, DART Quant Team}


\lead{The position's average price includes the 0.1\% fee.} The position held 3 shares at an average price of KRW~381,381 before the order, and 7 shares filled at an average of KRW~352,714.29. The final average price, KRW~361,561.2, equals $(3 \times 381{,}381 + 7 \times 352{,}714.29 \times 1.001)/10$ exactly, so \code{avgPrc} adds the fee to the fill price. The wrapper computes profit from \code{avgPrc} as the server reports it and does not add the fee a second time.

\lead{Placing, stopping and listing orders were recorded live; 3 other order paths were read from the client code only (\figref{orderapi}).} From 09:16 to 09:22 the browser's developer tools recorded the page's own requests while orders were placed and stopped by hand, and the remaining bodies were read from its JavaScript bundle. \code{AddOrder}, \code{CancelOrderBatch}, \code{ValidateProduct} and \code{Orders} were seen live. \code{CancelOrder}, \code{ModifyOrder} and \code{DeleteOrder} were not: the live order dialog stops orders through \code{CancelOrderBatch}, and the page shows modify and delete only for orders in its list of un-accepted or failed orders, so the bodies of these 3 paths are untested. The wrapper stops orders with \code{CancelOrderBatch}, joining the IDs with commas and leaving out orders already \code{Done}; it leaves the 3 untested paths out until they are tested live, so an amendment is a stop followed by a new order, which counts as 2 against the daily cap.

% OWNER: figs/orderapi.tex. Edit only this file.
\fig[wide]{orderapi}{\code{AddOrder}, \code{CancelOrderBatch} and \code{Orders} were captured live; 3 paths seen only in client code are untested}[Paths under \code{/api/Portfolio} (top) and the price and timing fields of an order body (bottom), captured 09:16 to 09:22 KST on 24 Nov 2026]{%
{\tabfont\begin{tabularx}{\linewidth}{>{\RaggedRight\arraybackslash}p{26mm} >{\RaggedRight\arraybackslash}p{9mm} >{\RaggedRight\arraybackslash}X >{\RaggedRight\arraybackslash}p{40mm} >{\RaggedRight\arraybackslash}p{27mm}}
\thd{Path} & \thd{Verb} & \thd{Parameters or body} & \thd{Reply} & \thd{Source}\\
\hl\code{AddOrder} & POST & \code{d}, \code{pfId}, \code{prodId}, \code{ls}, \code{ex}, \code{wei}, \code{exitAll}, \code{limitPrc}, \code{limitIdx}, \code{stopPrc}, \code{hm0}, \code{hm1} & 200, \texttt{\{"data":"","warnings":[]\}}; no order ID & Live (sell form: Section~\ref{sec:wrappertests})\\\rowsep
\hl\code{CancelOrderBatch} & POST & \code{d}, \code{pfId}, \code{ords}: order IDs joined by commas in one string, e.g.\ \code{"1650197"} & 200, empty body; the page labels it ``stop'' & Live\\\rowsep
\hl\code{Orders} & GET & \code{d0}, \code{d1}, \code{pfId} & 200, every order of the range, \code{Done} ones included & Live\\\rowsep
\code{ValidateProduct} & GET & \code{d}, \code{prodId}, \code{entry=true} & 200, \texttt{\{"data":"","warnings":[]\}} & Live\\\rowsep
\code{Session} & GET & \code{d}, \code{pfId}; called when the order dialog opens & 200, window 09:00 to 15:20 & Live\\\rowsep
\code{CancelOrder} & POST & \code{d}, \code{pfId}, \code{ordId} & Not observed & Client code, untested\\\rowsep
\code{ModifyOrder} & POST & \code{d}, \code{ordId}, \code{wei}, \code{exitAll}, \code{limitPrc}, \code{limitIdx}, \code{stopPrc}, \code{hm0}, \code{hm1}; no \code{pfId} & Not observed & Client code, untested\\\rowsep
\code{DeleteOrder} & POST & \code{d}, \code{pfId}, \code{ordId}; offered for un-accepted orders & Not observed & Client code, untested\\
\end{tabularx}}\par\vspace{5pt}
{\tabfont\begin{tabularx}{\linewidth}{>{\RaggedRight\arraybackslash}p{26mm} >{\RaggedRight\arraybackslash}p{52mm} >{\RaggedRight\arraybackslash}X}
\thd{Field} & \thd{Values} & \thd{Meaning}\\
\code{limitPrc} & price, with \code{limitIdx} = null & Limit order at a fixed price\\\rowsep
\code{limitIdx} & $n$ or $-n$, $n$ from 1 to 10 & $n$ ticks on the opposite side ($n$) or on the order's own side ($-n$)\\\rowsep
\code{stopPrc} & price & Stop order\\\rowsep
\code{hm0}, \code{hm1} & now and null; or start and end time & Immediate; or fills spread over the window from \code{hm0} to \code{hm1}\\
\end{tabularx}}}{Highlighted rows: the 3 paths the wrapper uses to place, stop and list orders; it also calls \code{ValidateProduct} before each order and \code{Session} for the trading window. The 3 paths at the bottom of the top table appear only in the JavaScript bundle; the wrapper does not call them. ``Live'' means the request and reply were recorded from the page; ``client code'' means the body was read from the JavaScript bundle. The wrapper's order calls take exactly one of \code{limitPrc}, \code{limitIdx} and \code{stopPrc} (Section~\ref{sec:wrappertests}).}{Web client code, DART Quant Team}


\lead{Every finished order ends in \code{Done}, so how it ended has to be read from \code{cnclT} and \code{cumQty}.} The bundle defines 5 states: \code{Generated}, \code{Accepted}, \code{Working}, \code{Done} and \code{Deleted}. The page disables modify and stop when \code{state} is \code{Done} or the stop time \code{cnclT} is not null. An \code{Orders} query at 09:31 showed all 3 stopped Samsung Electronics orders as \code{Done}, with \code{cnclT} set and 0 of 4 shares filled (\code{cumQty} 0 of \code{qty2ord} 4), the same state as the 2 filled orders. The wrapper therefore reads an order as finished when \code{state} is \code{Done}, as stopped when \code{cnclT} is also set, and as partly filled when \code{cnclT} is set and 0~$<$~\code{cumQty}~$<$~\code{qty2ord}. Two state sequences were observed: \code{Working} to \code{Done} for the morning orders, and \code{Accepted} to \code{Done} for order 1650719 at 10:33, which showed \code{Accepted} right after placement and was stopped at once. \code{Generated} and \code{Deleted} have not been observed.

\lead{The tests in this section used 5 new orders and 3 stops.} The 08:30 LIG Nex1 order came from a script and filled. Two Samsung Electronics orders were placed by hand: one of 0.1\% at 09:13 (\code{qty2ord} 4), which was stopped, and one of 0.3\% at 09:15 at 10 ticks on the opposite side, which filled 13 of 13 shares. Two more Samsung Electronics orders of 0.1\% (\code{qty2ord} 4 each) with a limit of KRW~151,800, 20\% below the market so that they could not fill, were placed at 09:18 (immediate) and 09:20 (window 15:00 to 15:20) to capture the requests, and both were stopped within a minute. The day's total against the cap is counted in Section~\ref{sec:wrappertests}.
\verdict{the wrapper confirms each \code{AddOrder} by finding the new record in the feed's \code{ords}, then in \code{Portfolio/Orders}, and raises \code{AmbiguousState} instead of resending when neither holds it (D11). It sends an order's \code{wei} as target minus current weight (D13), and reads \code{Done} as finished, stopped when \code{cnclT} is set, partly filled when also 0~$<$~\code{cumQty}~$<$~\code{qty2ord} (D15).}
